% ============================================================================
%  Section II — Related Work  (IOP MST)
%  Drop-in for Overleaf. Requires in the preamble:
%     \usepackage{booktabs}   % nicer tables
%     \usepackage{pifont}     % \ding checkmarks/crosses
%     \usepackage{multirow}
%  Citations resolve against refs.bib.
%  Convenience macros (place once in preamble; repeated here for portability):
\providecommand{\cmark}{\ding{51}}        % yes
\providecommand{\xmark}{\ding{55}}        % no
\providecommand{\pmark}{$\sim$}            % partial
% ============================================================================

\section{Related Work}
\label{sec:related}

\subsection{Vibration-Based Condition Monitoring of Induction Machines}
Vibration analysis is a mature and widely adopted modality for detecting
mechanical and electrical faults in induction machines, including bearing defects,
rotor asymmetries, and supply imbalances~\cite{gangsar2020signal,
hamani2025advancements}. Classical pipelines rely on industrial-grade
piezoelectric accelerometers sampled at tens to hundreds of kilohertz, which
resolve characteristic bearing defect frequencies---ball-pass frequency outer
race (BPFO), ball-pass frequency inner race (BPFI), and ball-spin frequency
(BSF)---and their harmonics through spectral and envelope analysis. This paradigm underpins the
public benchmarks that dominate the field: representative datasets are acquired
at $8$--$16$~kHz~\cite{lee2024multidomain}, $42$~kHz~\cite{sehri2024ottawa},
$51.2$~kHz~\cite{thuan2023hust}, and up to $200$~kHz~\cite{huang2018bearing}.
More recently, deep learning has become the default classifier family for such
signals~\cite{yang2021fault, zhang2020deep}. High-frequency acquisition,
however, entails costly instrumentation, intrusive cabling, and limited
portability, which restrict deployment to the comparatively small population of
critical machines and leave the many low-value machines in a plant unmonitored.

\subsection{Low-Cost MEMS and Smartphone-Based Vibration Sensing}
Micro-electro-mechanical-systems (MEMS) inertial sensors have emerged as an
inexpensive, compact alternative for vibration acquisition. Their metrological
limits relative to reference-grade accelerometers---narrower bandwidth and
higher noise floor---are well documented~\cite{grouios2022accelerometers}, and
condition-monitoring practitioners have generally regarded consumer-grade MEMS
bandwidth as insufficient for bearing diagnosis~\cite{pedotti2020lowcost}; the
metrological variability of such low-cost sensors and its effect on diagnostic
features has itself become a subject of measurement research~\cite{demilia2021mems}.
Nonetheless, a growing body of work demonstrates fault detection from
smartphone- or MEMS-acquired signals, using mobile applications and, more
recently, deep networks~\cite{naha2017mobile, ertargin2023deep,
ertargin2025automated}. In parallel, purpose-built low-cost MEMS sensor nodes
combine on-node processing with short-range or long-range radios for machine
diagnostics~\cite{pedotti2020lowcost, jakobsen2024lowcost}. The dataset employed
in this study~\cite{ertargin2026smartphone, ertargin2026dataset} extends this
line by providing long-duration, multi-condition smartphone-MEMS recordings at a
software-capped $100$~Hz. Crucially, the accompanying baseline reports
near-perfect classification accuracy under a random split of overlapping
windows, a protocol that is scrutinised in Section~\ref{subsec:leakage}.

\subsection{IoT, Edge, and TinyML for Condition Monitoring}
The Industrial Internet of Things reframes condition monitoring as a distributed
sensing problem in which many inexpensive nodes stream---or locally
distil---diagnostic information over constrained wireless
links~\cite{yousuf2024iot}. Because raw high-rate vibration cannot be sustained
over low-power wide-area radios, a large fraction of recent work pushes feature
extraction or inference onto the node itself. TinyML implementations of
lightweight classifiers on microcontrollers report bearing-fault accuracies with
millisecond-scale latency and sub-millijoule energy per
inference~\cite{gao2025edge, asutkar2023tinyml}, and low-cost MEMS nodes
integrate on-board spectral analysis with LoRa or Bluetooth Low Energy
telemetry~\cite{jakobsen2024lowcost}, and smart-sensor systems combine vibration
measurement with on-board bearing diagnosis~\cite{shukla2020smartsensor}. These
studies establish the feasibility of
edge diagnosis but typically assume industrial-bandwidth sensing and rarely
quantify the coupled trade-off between sampling/communication cost and diagnostic
accuracy that governs an ultra-low-rate node.

\subsection{Evaluation Methodology and Data Leakage}
\label{subsec:leakage}
A recurring threat to the credibility of learning-based fault diagnosis is data
leakage---the inadvertent sharing of information between training and test
partitions that inflates reported accuracy and does not transfer to
deployment~\cite{kapoor2023leakage}. In vibration diagnosis specifically, two
mechanisms dominate: (i) splitting \emph{overlapping} analysis windows drawn from
the same continuous recording, so that near-duplicate segments populate both
partitions; and (ii) condition- or recording-wise correlations that let a model
recognise \emph{which recording} a window originates from rather than the fault
itself. Recent studies show that enforcing component-wise (e.g., bearing-wise) or
recording-wise separation causes headline accuracies to fall
substantially~\cite{wheat2024impact, hendriks2022towards}. Despite these
warnings, random-split evaluation remains prevalent in smartphone- and
MEMS-based machine diagnosis, including the baseline accompanying the dataset used
here~\cite{ertargin2026smartphone}. A rigorous, deployment-honest re-evaluation
of low-rate MEMS condition monitoring is therefore still missing.

\subsection{Positioning and Contributions}
Table~\ref{tab:positioning} situates this work against representative prior art.
Existing studies advance one axis at a time. High-rate benchmarks and deep models
maximise accuracy but ignore cost and evaluation rigour; methodology papers expose
leakage but on industrial-grade data; smartphone-based studies embrace low-cost
sensing but retain optimistic random-split evaluation; and IoT/TinyML works
demonstrate edge deployment yet assume industrial bandwidth. To the best of the
authors' knowledge, no prior study jointly (i) evaluates ultra-low-rate
smartphone-MEMS machine diagnosis
under a leakage-free, cross-condition protocol, (ii) characterises the consumer
MEMS channel as a measurement instrument, and (iii) maps the accuracy--resource
trade-off into an IoT link budget and edge-model footprint. This paper closes that
gap, providing a reproducible and deployment-realistic baseline for
IoT-based condition monitoring of induction machines with low-cost MEMS sensing.

% ---------------------------------------------------------------------------
\begin{table*}[t]
\centering
\caption{Positioning of this work against representative prior studies.
\cmark~= addressed, \pmark~= partially addressed, \xmark~= not addressed.}
\label{tab:positioning}
\renewcommand{\arraystretch}{1.25}
\begin{tabular}{@{}p{3.9cm}p{2.6cm}ccccc@{}}
\toprule
\textbf{Representative work(s)} & \textbf{Sensor / rate} &
\textbf{Low} & \textbf{Leakage-free} & \textbf{Acc.--resource} &
\textbf{IoT/edge} & \textbf{Sensor} \\
 & & \textbf{cost} & \textbf{evaluation} & \textbf{trade-off} &
\textbf{analysis} & \textbf{charact.} \\
\midrule
High-rate vibration FD \& deep models~\cite{lee2024multidomain,sehri2024ottawa,thuan2023hust,zhang2020deep}
 & Industrial accel., $8$--$200$\,kHz & \xmark & \xmark & \xmark & \xmark & \xmark \\
Evaluation / leakage studies~\cite{kapoor2023leakage,wheat2024impact,hendriks2022towards}
 & Industrial accel., k\,Hz & \xmark & \cmark & \xmark & \xmark & \xmark \\
Smartphone / mobile machine FD~\cite{naha2017mobile,ertargin2023deep,ertargin2025automated}
 & Smartphone MEMS, $\leq$\,k\,Hz & \cmark & \xmark & \xmark & \pmark & \xmark \\
Source-dataset baseline~\cite{ertargin2026smartphone}
 & Smartphone MEMS, $100$\,Hz & \cmark & \xmark & \xmark & \xmark & \pmark \\
Low-cost MEMS sensor nodes~\cite{pedotti2020lowcost,jakobsen2024lowcost}
 & MEMS node, k\,Hz & \cmark & \xmark & \pmark & \cmark & \pmark \\
TinyML edge fault diagnosis~\cite{gao2025edge,asutkar2023tinyml}
 & (Lab) accel., k\,Hz & \cmark & \xmark & \pmark & \cmark & \xmark \\
\midrule
\textbf{This work} & \textbf{Smartphone-grade MEMS, $100$\,Hz}
 & \cmark & \cmark & \cmark & \cmark & \cmark \\
\bottomrule
\end{tabular}
\end{table*}
